%%%%%%%%%%%%%%%%%%%%%%%%%%%%%%%%%%%%%%%%%%%%%%%%%%%%%%%%%%%%%%%%%%%%%%%%%%%%%%
%
% 09_dynamics.tex
%
% SECTION 9
% Dynamics: Growth Operators and the Algebra D
%
%%%%%%%%%%%%%%%%%%%%%%%%%%%%%%%%%%%%%%%%%%%%%%%%%%%%%%%%%%%%%%%%%%%%%%%%%%%%%%

\begin{mosbox}{9. Dynamics: Growth Operators and the Operator Algebra}

\BodyFont

Unlike a static database, the Mathematical Operating System evolves by
modifying its underlying mathematical structure.

Learning is interpreted as a change in shape: concepts are reorganized,
functional coalitions strengthen or decay, sheaf data adapts, and
successful procedures become permanent mathematical operators.

The evolution of MOS is therefore governed not only by changes in memory,
but also by changes in the collection of admissible operations.

\end{mosbox}

%%%%%%%%%%%%%%%%%%%%%%%%%%%%%%%%%%%%%%%%%%%%%%%%%%%%%%%%%%%%%%%%%%%%%%%%%%%%%%
%% HERO EQUATION
%%%%%%%%%%%%%%%%%%%%%%%%%%%%%%%%%%%%%%%%%%%%%%%%%%%%%%%%%%%%%%%%%%%%%%%%%%%%%%

\begin{eqbox}

\[
\Huge
\boxed{
\mathcal D(t)
\longrightarrow
\mathcal D(t+1)
}
\]

\vspace{3mm}

\BodyFont

The operator algebra is itself dynamic.
Learning expands not only knowledge but also the mathematical operations
available for future reasoning.

\end{eqbox}

%%%%%%%%%%%%%%%%%%%%%%%%%%%%%%%%%%%%%%%%%%%%%%%%%%%%%%%%%%%%%%%%%%%%%%%%%%%%%%
%% FOUR AXES
%%%%%%%%%%%%%%%%%%%%%%%%%%%%%%%%%%%%%%%%%%%%%%%%%%%%%%%%%%%%%%%%%%%%%%%%%%%%%%

\begin{mosbox}{The Four Growth Axes}

\BodyFont

Evolution proceeds simultaneously along four independent directions.

\[
\boxed{
\delta C_{\rm fine},
\qquad
\delta C_{\rm coarse},
\qquad
\delta\mathcal F,
\qquad
\delta\mathcal D
}
\]

\vspace{4mm}

\begin{center}

\begin{tabular}{p{4cm}p{10.8cm}}

$\delta C_{\rm fine}$

&

Growth of concepts (semantic memory).

\\[3mm]

$\delta C_{\rm coarse}$

&

Formation and restructuring of coalitions
(functional connectivity).

\\[3mm]

$\delta\mathcal F$

&

Evolution of stalks and restriction maps within the cellular sheaf.

\\[3mm]

$\delta\mathcal D$

&

Growth of the operator algebra through the acquisition of new
procedural skills.

\end{tabular}

\end{center}

The fourth axis distinguishes MOS from a fixed mathematical model:
the collection of admissible operators evolves over time.

\end{mosbox}

%%%%%%%%%%%%%%%%%%%%%%%%%%%%%%%%%%%%%%%%%%%%%%%%%%%%%%%%%%%%%%%%%%%%%%%%%%%%%%
%% FRACTAL GROWTH TREE
%%%%%%%%%%%%%%%%%%%%%%%%%%%%%%%%%%%%%%%%%%%%%%%%%%%%%%%%%%%%%%%%%%%%%%%%%%%%%%

\begin{center}

\begin{tikzpicture}[
  organic/.style={draw=MOSRed, line width=2pt, out=90, in=-90},
  neuron/.style={circle, fill=MOSRed, draw=MOSGold, line width=1.5pt, minimum size=10pt, inner sep=0pt},
  textnode/.style={font=\Large, text=white}
]

% Nodes
\node[textnode] (M) at (0,0) {$M(t)$};
\node[textnode] (ops) at (0,1.5) {bind\quad split\quad merge};
\node[textnode] (dag) at (0,3) {Composite DAG};
\node[textnode] (micro) at (0,4.5) {Microneuron};
\node[textnode] (verify) at (0,6) {\texttt{VerifyOp}};

\node[neuron] (nL) at (-2, 7.5) {};
\node[neuron] (nC) at (0, 7.5) {};
\node[neuron] (nR) at (2, 7.5) {};

\node[neuron] (top) at (0, 9) {};
\node[textnode, text=MOSGold, font=\Huge\bfseries] (D) at (0, 10) {$\mathcal D(t+1)$};

% Organic connections
\draw[organic] (M) to (ops);
\draw[organic] (ops) to (dag);
\draw[organic] (dag) to (micro);
\draw[organic] (micro) to (verify);

\draw[organic] (verify) to (nL);
\draw[organic] (verify) to (nC);
\draw[organic] (verify) to (nR);

\draw[organic] (nL) to (top);
\draw[organic] (nC) to (top);
\draw[organic] (nR) to (top);

\draw[organic] (top) to (D);

\end{tikzpicture}

\end{center}

\begin{center}

\Large

\textcolor{MOSGold}{
Learning as Structural Growth
}

\end{center}

%%%%%%%%%%%%%%%%%%%%%%%%%%%%%%%%%%%%%%%%%%%%%%%%%%%%%%%%%%%%%%%%%%%%%%%%%%%%%%
%% STRUCTURAL OPERATORS
%%%%%%%%%%%%%%%%%%%%%%%%%%%%%%%%%%%%%%%%%%%%%%%%%%%%%%%%%%%%%%%%%%%%%%%%%%%%%%

\begin{highlightbox}

\SectionTitle{Structural Operators}

\BodyFont

The elementary structural moves acting on the simplicial complex are

\[
\boxed{
\{
\mathrm{bind},
\mathrm{collapse},
\mathrm{split},
\mathrm{merge}
\}.
}
\]

Every operation is recorded in an immutable mutation history.

\vspace{2mm}

\textbf{Remark.}

The earlier terminology \emph{Pachner move} was intentionally discarded.
A knowledge complex is not a manifold, therefore Pachner invariance does
not apply.

\end{highlightbox}

%%%%%%%%%%%%%%%%%%%%%%%%%%%%%%%%%%%%%%%%%%%%%%%%%%%%%%%%%%%%%%%%%%%%%%%%%%%%%%
%% HEBBIAN
%%%%%%%%%%%%%%%%%%%%%%%%%%%%%%%%%%%%%%%%%%%%%%%%%%%%%%%%%%%%%%%%%%%%%%%%%%%%%%

\begin{mosbox}{Hebbian Coupling and Structural Decay}

\BodyFont

Each coalition

\[
\sigma
\]

is assigned a weight

\[
w(\sigma,t)\in[0,1].
\]

Repeated co-activation strengthens the coalition

\[
\boxed{
\text{Fire Together}
\Longrightarrow
\text{Wire Together}.
}
\]

Idle coalitions decay according to

\[
\boxed{
w(\sigma,t+1)
=
(1-\gamma)\,
w(\sigma,t).
}
\]

Coalitions whose weights fall below the binding threshold collapse,
providing an intrinsic form of structural garbage collection.

Pairs whose weights exceed the threshold become
1-simplices (edges) of the knowledge complex.

\end{mosbox}

%%%%%%%%%%%%%%%%%%%%%%%%%%%%%%%%%%%%%%%%%%%%%%%%%%%%%%%%%%%%%%%%%%%%%%%%%%%%%%
%% CONSOLIDATION
%%%%%%%%%%%%%%%%%%%%%%%%%%%%%%%%%%%%%%%%%%%%%%%%%%%%%%%%%%%%%%%%%%%%%%%%%%%%%%

\begin{mosbox}{Persistence and Microneurons}

\BodyFont

A weight filtration is analyzed using persistent homology.

Coalitions and concept-cycles that persist across the filtration are
promoted into permanent mathematical structure.

A particularly important promoted object is the
\textbf{microneuron}:

\begin{itemize}

\item an operadic sub-tree,

\item represented internally as a composite computational DAG,

\item repeatedly resolves conflict,

\item repeatedly passes verification,

\item promoted into a new atomic generator of
$\mathcal D$.

\end{itemize}

The promotion gate is

\[
\boxed{
\texttt{VerifyOp}
}
\]

Only externally verified composite procedures become permanent generators
of the operator algebra.

\end{mosbox}

%%%%%%%%%%%%%%%%%%%%%%%%%%%%%%%%%%%%%%%%%%%%%%%%%%%%%%%%%%%%%%%%%%%%%%%%%%%%%%
%% TAKEAWAY
%%%%%%%%%%%%%%%%%%%%%%%%%%%%%%%%%%%%%%%%%%%%%%%%%%%%%%%%%%%%%%%%%%%%%%%%%%%%%%

\begin{remarkbox}

\BodyFont

\textbf{Mathematical Interpretation}

MOS evolves simultaneously on four interconnected levels:

\[
\boxed{
\text{Concepts}
\rightarrow
\text{Coalitions}
\rightarrow
\text{Sheaf Data}
\rightarrow
\text{Operator Algebra}.
}
\]

Unlike the preceding three growth mechanisms, the final axis changes the
very collection of admissible mathematical operations, allowing the
system to acquire new procedural primitives over time.

\end{remarkbox}